\par\noindent\begin{minipage}{\linewidth}
\small
\setlength{\tabcolsep}{3pt}
\renewcommand{\arraystretch}{1.15}
\captionof{table}{Oposición persistente con cobertura completa de medidas alternativas.}\label{tab:S15}
\begin{tabular*}{\linewidth}{@{\extracolsep{\fill}}lrrrrr@{}}
\toprule
\textbf{Modelos} & \textbf{N} & \textbf{Apertura} & \textbf{\shortstack{Apertura +\\alternativas}} & \textbf{PR} & \textbf{\shortstack{PR +\\alternativas}} \\
\midrule
Los diez & 10 & 262 & 225 & 221 & 190 \\[2pt]
Seis de similitud & 6 & 231 & 190 & 160 & 115 \\[2pt]
Sin SPECTER & 9 & 234 & 182 & 212 & 174 \\[2pt]
Sin SPECTER2 & 9 & 251 & 205 & 211 & 173 \\[2pt]
Sin SciNCL & 9 & 259 & 217 & 215 & 182 \\[2pt]
Sin SciBERT & 9 & 258 & 222 & 218 & 187 \\[2pt]
Sin BERT & 9 & 261 & 223 & 211 & 176 \\[2pt]
Sin MPNet & 9 & 262 & 223 & 218 & 186 \\[2pt]
Sin MiniLM & 9 & 259 & 223 & 221 & 188 \\[2pt]
Sin PubMedBERT & 9 & 261 & 220 & 206 & 187 \\[2pt]
Sin BioBERT & 9 & 259 & 217 & 212 & 169 \\[2pt]
Sin SimCSE & 9 & 259 & 217 & 217 & 183 \\[2pt]
\bottomrule
\end{tabular*}
\par\smallskip{\small Recuentos sobre 325 con magnitud mínima cero. Todas las alternativas angulares y de dimensión cubren ahora las 26 selecciones. Los mismos modelos opuestos deben conservar sus direcciones; no se sustituyen por otros. Los empates de D80 quedan sin resolver. Los seis modelos son SPECTER, SPECTER2, SciNCL, MPNet, MiniLM y SimCSE. Los CSV conservan cortes y direcciones completos.\par}
\end{minipage}\par
